\begin{frame}{프로젝트 구조: 환경 · 2개 이상 알고리즘 · 필수 요구사항}
  \begin{itemize}
    \item \kb{환경}: Gymnasium의 \textbf{표 기반(이산 상태·행동)} 환경 하나 ---
      \texttt{FrozenLake-v1}, \texttt{CliffWalking-v1}, \texttt{Taxi-v4} 등.
    \item \kb{적용}: Ch\,4$\sim$7 알고리즘 \textbf{최소 두 가지}를 같은 환경에
      적용 (예: Q-learning vs SARSA, MC 제어 vs n-step TD) $\to$ 학습 곡선과
      최종 정책 비교.
    \item \kb{필수 --- 수식적 정당화 최소 1개}:
    \begin{itemize}
      \item Q-learning/SARSA 비교 $\to$ Ch\,6.2$\sim$6.3 on/off-policy로
        정책이 왜 다르게 학습됐는지 설명.
      \item $\varepsilon$/$\alpha$ 변경 실험 $\to$ Ch\,6.3 수렴 조건과 연결.
      \item Dyna-Q 사용 $\to$ Ch\,7.3 계획 스텝 수와 학습 속도를 \textbf{직접 측정}.
    \end{itemize}
  \end{itemize}
\end{frame}
